Throughout the proof, the phrase ``on the good-Gram event'' means multiplication by the indicator
\[
\mathbf 1_{G_h},
\qquad
G_h=\{\lambda_{\min}(\widehat D_h)\ge\lambda_*\},
\]
which is the truncation used by the estimator in \(\mathrm{def:beta\text{-}free\text{-}handle}\). Condition on \(I_\mu\) and \(I_{\mathrm{aux}}\). Put \(N=|I_{\mathrm{lp}}|\), \(\bar f=P_X\widehat\pi_0\), \(\bar m=P_X\widehat\mu_0\), and
\[
\bar\xi_i=
\frac{\bar f}{\widehat\pi_0(X_i)}\{Y_i-\widehat\mu_0(X_i)\}+\bar m.
\]
The construction gives \(|\widehat\mu_0|\le M_Y\) and \(c/2\le\widehat\pi_0\le2/c\), so \(|\bar\xi_i|\le C(M_Y,c)\). Define
\[
D_h=\mathbb E_P[K_h(A)S_h(A)S_h(A)^\top],
\qquad
\gamma_h=D_h^{-1}\mathbb E_P[K_h(A)S_h(A)\bar\xi].
\]
The population-Gram conclusion of \(\mathrm{lem:design\text{-}adaptive\text{-}local\text{-}polynomial\text{-}bias}\), together with the two-sided marginal-density bound in \(\mathrm{lem:marginal\text{-}treatment\text{-}positivity}\), gives \(\|D_h^{-1}\|_{\mathrm{op}}+\|\gamma_h\|\le C\). On \(G_h\), exact normal-equation algebra gives
\[
\widehat\gamma_h-\gamma_h
=
\widehat D_h^{-1}\frac1N\sum_{i\in I_{\mathrm{lp}}}
K_h(A_i)S_h(A_i)\{\bar\xi_i-S_h(A_i)^\top\gamma_h\}.
\]
The vector summands are conditionally iid and centered. Their second moments are bounded by \(C/h\), using the upper bound on the marginal density of \(A\). Since \(\|\widehat D_h^{-1}\|_{\mathrm{op}}\le1/\lambda_*\) on \(G_h\),
\[
\mathbb E_P\!\left[
\mathbf1_{G_h}|e_0^\top(\widehat\gamma_h-\gamma_h)|^2
\,\middle|\,I_\mu,I_{\mathrm{aux}}
\right]
\le\frac{C}{Nh}\le\frac{C}{nh}.
\]
This proves the first assertion.

For the empirical marginalizations, set
\[
\Delta_m=P_{I_X}\widehat\mu_0-P_X\widehat\mu_0,
\qquad
\Delta_f=P_{I_X}\widehat\pi_0-P_X\widehat\pi_0,
\qquad
q_i=\frac{Y_i-\widehat\mu_0(X_i)}{\widehat\pi_0(X_i)}.
\]
Conditional independence and boundedness give
\[
\mathbb E_P(\Delta_m^2+\Delta_f^2\mid I_\mu,I_{\mathrm{aux}})\le C/n.
\]
Because \(S_h(a)^\top e_0=1\),
\[
P_{I_{\mathrm{lp}}}(K_hS_h)=\widehat D_he_0,
\]
and hence, on \(G_h\),
\[
e_0^\top(\widehat\gamma_h^{\mathrm{emp}}-\widehat\gamma_h^{\mathrm{pop}})
=
\Delta_m+\Delta_fT_h,
\qquad
T_h=e_0^\top\widehat D_h^{-1}P_{I_{\mathrm{lp}}}(K_hS_hq).
\]
It remains to bound the truncated second moment of \(T_h\). Let
\[
J=\sum_{i\in I_{\mathrm{lp}}}\mathbf1\{K_h(A_i)>0\}.
\]
If \(p=\lfloor\alpha\rfloor=0\), \(T_h\) is a nonnegative-kernel weighted average of the bounded \(q_i\), and so \(\mathbf1_{G_h}|T_h|\le\|q\|_\infty\). If \(p\ge1\), put \(m=p+1\ge2\). On \(G_h\), the rank of \(\widehat D_h\) is \(m\), so \(J\ge m\), and
\[
\mathbf1_{G_h}|T_h|
\le C\mathbf1_{\{J\ge m\}}\frac{J}{Nh}.
\]
When \(Nh\le1\), the upper marginal-density bound makes \(J\) binomial with mean at most \(CNh\). Since \(J^2\le2J(J-1)\) on \(\{J\ge2\}\),
\[
\mathbb E_P[\mathbf1_{G_h}T_h^2\mid I_\mu,I_{\mathrm{aux}}]
\le\frac{C}{(Nh)^2}\mathbb E[J(J-1)]\le C.
\]
When \(Nh>1\), define
\[
\tau_h=D_h^{-1}\mathbb E_P[K_h(A)S_h(A)q].
\]
The same centered-sum calculation used above gives
\[
\mathbb E_P[\mathbf1_{G_h}|T_h-e_0^\top\tau_h|^2\mid I_\mu,I_{\mathrm{aux}}]
\le\frac{C}{Nh},
\]
while \(\|\tau_h\|\le C\). Thus the truncated second moment of \(T_h\) is uniformly bounded in both regimes. Finally, \(I_X\) is independent of \(I_{\mathrm{lp}}\), so
\[
\mathbb E_P[\mathbf1_{G_h}|e_0^\top(\widehat\gamma_h^{\mathrm{emp}}-
\widehat\gamma_h^{\mathrm{pop}})|^2]
\le C/n.
\]
This proves the second assertion. Notice that this conditional calculation uses the definition and inverse bound of the good-Gram event, but not the separate probability estimate for its complement.